\SourceLeaf{260}
\clearpage
\Article{9}
All dies war bereits in den Disquis.\ Ar.\ auf nahezu dieselbe Weise behandelt worden; einiges war dort jedoch ohne Beweis hinzugefügt worden, den wir hier nachzutragen für zweckmäßig halten. Wir haben dort angekündigt, dass die Bestimmung des Wertes der Wurzelgröße ${}^{\beta}\!\surd\theta'$, die, da $\theta'$ eine imaginäre Größe ist, die Teilung sowohl eines Verhältnisses als auch eines Winkels in $\beta$ Teile zu erfordern scheint, nur von der letzteren abhängt und dass die erstere stets auf das Ziehen einer einzigen Quadratwurzel zurückgeführt werden kann. Dies beweisen wir wie folgt.

Bezeichnen wir wie oben die imaginäre Größe $\surd-1$ mit $i$ und schreiben $\theta'=P+iQ$ sowie ${}^{\beta}\!\surd\theta'=p+iq$, wobei für die Wurzelgröße ein beliebiger Wert gewählt ist und $P,Q,p,q$ reell sind, so ist bekannt: Werden die positiven Größen $E,e$ und die Winkel $F,f$ so bestimmt, dass $P=E\cos F$, $Q=E\sin F$, $p=e\cos f$, $q=e\sin f$ ist, so gilt $e={}^{\beta}\!\surd E$, und $f$ ist gleich einem der Winkel
\[
\frac1\beta F,\quad\frac1\beta(F+360^0),\quad\frac1\beta(F+720^0)\dotsfour\frac1\beta(F+(\beta-1)360^0)
\]
Somit wird $f$\textsuperscript{.}\ ger durch Teilung des Winkels $F$ in $\beta$ Teile bestimmt; das Ziehen der Wurzel ${}^{\beta}\!\surd E$ können wir jedoch auf folgende Weise vermeiden. Jedes Produkt $r^kR^K$ hat denselben reellen Teil wie $r^{-k}R^{-K}$, während die imaginären, den Faktor $i$ enthaltenden Teile in diesen Produkten gleich, aber entgegengesetzt sind. Daraus folgt unmittelbar $[r^{-1},R^{-1}]=p-iq=e(\cos f-i\sin f)$ und daher
\[
[r,R]\times[r^{-1},R^{-1}]=e^2
\]

Nach dem Satz aus Artikel~4 wird jenes Produkt jedoch
\[
\begin{aligned}
&{}=[1,1]+R[r^{g^\alpha-1},1]+R^2[r^{g^{2\alpha}-1},1]+\text{usw.}\\
&\qquad{}+R^{\beta\gamma-1}[r^{g^{\alpha\beta\gamma-\alpha}-1},1]\\
&{}=\beta\gamma+R(\beta\gamma,g^\alpha-1)+R^2(\beta\gamma,g^{2\alpha}-1)+\text{usw.}\\
&\qquad{}+R^{\beta\gamma-1}(\beta\gamma,g^{\alpha\beta\gamma-\alpha}-1)
\end{aligned}
\]
Diese Größe lässt sich bestimmen, wenn $R$ und sämtliche Perioden mit $\beta\gamma$ Gliedern als bekannt vorausgesetzt werden. Die Bestimmung von $e$ verlangt somit nur das Ziehen einer Quadratwurzel.

Im besonderen Fall $\alpha=1$ sind die einzelnen Perioden $(\beta\gamma,g^\alpha-1)$, $(\beta\gamma,g^{2\alpha}-1)$ usw.\ offenbar $=r+r^2+r^3+r^4+\text{usw.}+r^{n-1}$, und daher
\[
\begin{aligned}
ee&=\beta\gamma+(R+R^2+R^3+\text{usw.}+R^{\beta\gamma-1})(r+r^2+r^3+\text{usw.}+r^{n-1})\\
&=\beta\gamma+1=n
\end{aligned}
\]
\SourceLeaf{261}
sofern $r$ und $R$ als von $1$ verschiedene Wurzeln vorausgesetzt werden; folglich gilt stets $e=\surd n$ (Disq. arithm., Artikel~360 am Ende).

\Article{10}
Bisher haben wir unsere Untersuchung in größter Allgemeinheit geführt, sodass sie beliebige Werte der Zahlen $\alpha,\beta,\gamma$ umfasst; von nun an gehen wir jedoch zu dem eingeschränkteren Fall $\alpha=1$ über, der uns den Weg zu überaus fruchtbaren und eleganten Untersuchungen eröffnen wird. Das Zeichen $[r,R]$ bezeichne also die Funktion
\[
r+Rr^g+R^2r^{g^2}+R^3r^{g^3}+\text{usw.}+R^{n-2}r^{g^{n-2}}
\]
Dabei ist $n$ eine Primzahl, $r$ eine beliebige Wurzel der Gleichung $x^n-1=0$ (die Wurzel $1$ nicht ausgenommen), $R$ eine beliebige Wurzel der Gleichung $x^\beta-1=0$, wobei $\beta$ einen gegebenen Teiler von $n-1$ bezeichnet, und schließlich $g$ eine ganze Zahl, die eine fest gewählte primitive Wurzel zum Modul $n$ ist. Ferner schreiben wir der Kürze halber
\[
\begin{aligned}
1+r+r^2+r^3+\text{usw.}+r^{n-1}&=s\\
1+R+R^2+R^3+\text{usw.}+R^{n-2}&=S
\end{aligned}
\]
Daraus folgt offenbar, dass $s$ den Wert $=n$ für $r=1$ annimmt, dagegen $s=0$ für jeden anderen Wert von $r$; ebenso gilt $S=n-1$ für $R=1$, dagegen $S=0$ für jeden anderen Wert von $R$.

Nach Artikel~3 haben wir daher $[1,R]=S$, $[r,1]=s-1$; ferner gilt für jeden ganzzahligen Wert von $m$, der nicht durch $n$ teilbar ist, $[r^m,R]=R^{-\operatorname{ind}m}[r,R]$ oder allgemeiner $[r^m,R^M]=R^{-M\operatorname{ind}m}[r,R^M]$, wobei $\operatorname{ind}m$ der Exponent der Potenz der Zahl $g$ ist, die zum Modul $n$ der Zahl $m$ kongruent ist. Wenden wir diese Transformation auf das an, was wir in Artikel~5 gezeigt haben, so folgt, dass sich das Produkt zweier oder mehrerer solcher Funktionen $[r^h,R^H]$ auf die Form
\[
A[1,R^\lambda]+B[r,R^\lambda]
\]
bringen lässt. Dabei sind $A$ und $B$ rationale Funktionen von $R$ mit ganzzahligen Koeffizienten, und $\lambda$ ist die Summe aller Werte von $H$. Es wird von großer Bedeutung sein, solche Transformationen auf einen handlichen Algorithmus zurückzuführen; dazu müssen wir vor allem die Beschaffenheit des Produkts zweier Funktionen näher betrachten.

\Needspace{21\baselineskip}
\Article{11}
Das Produkt $[r,R^\mu]\times[r,R^\nu]$ wird nach dem Satz aus Artikel~4 $=$
\SourceLeaf{262}
\[
\begin{aligned}
&[r^2,R^{\mu+\nu}]+R^\mu[r^{g+1},R^{\mu+\nu}]+R^{2\mu}[r^{g^2+1},R^{\mu+\nu}]\\
&\qquad+R^{3\mu}[r^{g^3+1},R^{\mu+\nu}]+\text{usw.}+R^{(n-2)\mu}[r^{g^{n-2}+1},R^{\mu+\nu}]
\end{aligned}
\]
Unter den $n-1$ Exponenten $2,g+1,g^2+1,g^3+1$ usw.\ $g^{n-2}+1$ ist nur einer durch $n$ teilbar, nämlich $g^{\frac12(n-1)}+1$; das entsprechende Glied unserer Summe ist daher $R^{\frac12(n-1)\mu}[1,R^{\mu+\nu}]$. Dieses Glied ist $=0$, sofern nicht $R^{\mu+\nu}=1$ ist, und $=(n-1)R^{\frac12(n-1)\mu}=\pm(n-1)$ für $R^{\mu+\nu}=1$. Die übrigen Teile unserer Summe, deren Summe wir $=\Omega$ setzen, werden auf folgende Weise umgeformt:
\[
\begin{aligned}
\relax[r^2,R^{\mu+\nu}]&=R^{-(\mu+\nu)\operatorname{ind}2}[r,R^{\mu+\nu}]\\
R^\mu[r^{g+1},R^{\mu+\nu}]&=R^{\mu-(\mu+\nu)\operatorname{ind}(g+1)}[r,R^{\mu+\nu}]\\
R^{2\mu}[r^{g^2+1},R^{\mu+\nu}]&=R^{2\mu-(\mu+\nu)\operatorname{ind}(g^2+1)}[r,R^{\mu+\nu}]\\
R^{3\mu}[r^{g^3+1},R^{\mu+\nu}]&=R^{3\mu-(\mu+\nu)\operatorname{ind}(g^3+1)}[r.R^{\mu+\nu}]\\
&\text{usw.}
\end{aligned}
\]

Daraus erhalten wir
\begingroup\postdisplaypenalty=10000
\[
\text{I.}\qquad\Omega=[r,R^{\mu+\nu}]\times\Sigma R^{\mu\operatorname{ind}x-(\mu+\nu)\operatorname{ind}(x+1)}
\]
\endgroup
wenn für $x$ der Reihe nach die Werte $1,g,g^2,g^3\dotsfour g^{n-2}$ mit Ausnahme von $g^{\frac12(n-1)}$ eingesetzt werden oder, was offenbar auf dasselbe hinausläuft, wenn wir für $x$ die Werte $1,2,3,4\dotsfour n-2$ einsetzen, da diese Werte jenen zum Modul $n$, wenn auch in anderer Reihenfolge, kongruent sind.

Nehmen wir die ganze Zahl $y$ zu $x$ invers zum Modul $n$, also so bestimmt, dass $xy\equiv1\pmod n$ gilt, so ist $\operatorname{ind}x\equiv-\operatorname{ind}y\pmod{n-1}$ und $\operatorname{ind}(x+1)+\operatorname{ind}y\equiv\operatorname{ind}(xy+y)\equiv\operatorname{ind}(1+y)\pmod{n-1}$; daraus ergibt sich
\[
\begin{aligned}
\mu\operatorname{ind}x-(\mu+\nu)\operatorname{ind}(x+1)&\equiv-\mu\operatorname{ind}y-(\mu+\nu)\{\operatorname{ind}(y+1)-\operatorname{ind}y\}\\
&\equiv\nu\operatorname{ind}y-(\mu+\nu)\operatorname{ind}(y+1)
\end{aligned}
\]
Da nun die zu $1,2,3\dots n-2$ inversen Zahlen mit diesen selbst übereinstimmen und nur ihre Reihenfolge geändert ist, gilt auch
\[
\text{II.}\qquad\Omega=[r,R^{\mu+\nu}]\times\Sigma R^{\nu\operatorname{ind}y-(\mu+\nu)\operatorname{ind}(y+1)}
\]
wobei für $y$ der Reihe nach die Zahlen $1,2,3\dots n-2$ eingesetzt werden. Dieselbe Formel folgt unmittelbar aus I, da sich die Zahlen $\mu,\nu$ offenbar miteinander vertauschen lassen.

Nehmen wir schließlich die ganze Zahl $z$ zu $x+1$ invers zum
\SourceLeaf{263}
Modul $n$, also $xz+z\equiv1\pmod n$, so gilt $\operatorname{ind}(1-z)\equiv\operatorname{ind}x+\operatorname{ind}z\pmod{n-1}$, $\operatorname{ind}(x+1)\equiv-\operatorname{ind}z\pmod{n-1}$ und daher
\[
\begin{aligned}
\mu\operatorname{ind}x-(\mu+\nu)\operatorname{ind}(x+1)&\equiv\mu(\operatorname{ind}(1-z)-\operatorname{ind}z)+(\mu+\nu)\operatorname{ind}z\\
&\equiv\mu\operatorname{ind}(1-z)+\nu\operatorname{ind}z
\end{aligned}
\]
Da folglich, wenn $x$ die Werte $1,2,3\dots n-2$ durchläuft, die Zahl $z$ die Werte $2,3,4\dots n-1$, wenn auch in anderer Reihenfolge, durchlaufen muss, erhalten wir den dritten Ausdruck
\begingroup\postdisplaypenalty=10000
\[
\text{III.}\qquad\Omega=[r,R^{\mu+\nu}]\times\Sigma R^{\mu\operatorname{ind}(1-z)+\nu\operatorname{ind}z}
\]
\endgroup
wobei für $z$ der Reihe nach die Werte $2,3,4\dots n-1$ eingesetzt werden, oder, wenn du es vorziehst,
\[
\begin{aligned}
\text{IV.}\qquad\Omega&=[r,R^{\mu+\nu}]\times\Sigma R^{\mu\operatorname{ind}(n+1-z)+\nu\operatorname{ind}z}\\
&=[r,R^{\mu+\nu}]\times\Sigma R^{\mu\operatorname{ind}z+\nu\operatorname{ind}(n+1-z)}
\end{aligned}
\]
Da $\operatorname{ind}(1-z)=\frac12(n-1)+\operatorname{ind}(z-1)$ gilt, lässt sich unser Produkt auch so darstellen:
\[
\begin{aligned}
\relax[r,R^\mu]\times[r,R^\nu]&=R^{\frac12(n-1)\mu}\{[1,R^{\mu+\nu}]+[r,R^{\mu+\nu}]\times\Sigma R^{\mu\operatorname{ind}(z-1)+\nu\operatorname{ind}z}\}\\
&=R^{\frac12(n-1)\nu}\{[1,R^{\mu+\nu}]+[r,R^{\mu+\nu}]\times\Sigma R^{\mu\operatorname{ind}z+\nu\operatorname{ind}(z-1)}\}
\end{aligned}
\]
Dabei sind für $z$ stets die Werte $2,3,4\dotsfour n-1$ einzusetzen.

Ferner können in allen diesen Formeln die Zahlen
\[
\mu\operatorname{ind}x-(\mu+\nu)\operatorname{ind}(x+1),\quad\nu\operatorname{ind}y-(\mu+\nu)\operatorname{ind}(y+1),\quad\mu\operatorname{ind}(1-z)+\nu\operatorname{ind}z
\]
usw.\ offenbar durch ihre kleinsten Reste zum Modul $\beta$ ersetzt werden.

Ist $\mu+\nu\equiv0\pmod\beta$, so gilt
\[
\begin{aligned}
\relax[r,R^\mu][r,R^\nu]&=(n-1)R^{\frac12(n-1)\mu}\\
&\quad+(r+r^2+r^3+\dots+r^{n-1})\\
&\qquad\times(1+R^\mu+R^{2\mu}+R^{3\mu}+\dots+R^{(n-2)\mu}-R^{\frac12(n-1)\mu})
\end{aligned}
\]

\clearpage
\Article{12}
Das Produkt $[1,R^\mu]\times[r,R^\nu]$ wird nach dem Satz aus Artikel~4
\[
\begin{aligned}
&=[r,R^{\mu+\nu}]+R^\mu[r,R^{\mu+\nu}]+R^{2\mu}[r,R^{\mu+\nu}]+\text{usw.}+R^{(n-2)\mu}[r,R^{\mu+\nu}]\\
&=[r,R^{\mu+\nu}]\times(1+R^\mu+R^{2\mu}+R^{3\mu}+\text{usw.}+R^{(n-2)\mu})\\
&=[r,R^{\mu+\nu}]\times\frac{n-1}\beta(1+R^\mu+R^{2\mu}+R^{3\mu}+\text{usw.}+R^{(\beta-1)\mu})
\end{aligned}
\]

\SourceLeaf{264}
Daraus entwickelt sich das Produkt $[1,R^\mu]\times[1,R^\nu]$ zu
\[
\frac{n-1}\beta[1,R^{\mu+\nu}]\times(1+R^\mu+R^{2\mu}+R^{3\mu}+\text{usw.}+R^{(\beta-1)\mu})
\]
Nun lässt sich das Produkt $[r^m,R^\mu]\pdot[r^{m'},R^{\mu'}]$ allgemein ohne Schwierigkeit ermitteln. Denn es ist $[r^m,R^\mu]=R^{-\mu\operatorname{ind}m}[r,R^\mu]$ für einen Wert von $m$, der nicht durch $n$ teilbar ist, und $=[1,R^\mu]$ für einen teilbaren Wert; und da eine ähnliche Transformation für den anderen Faktor $[n^{m'},R^{\mu'}]$ gilt, führt die Multiplikation entweder auf die Aufgabe des vorigen Artikels oder auf die Fälle zurück, die wir in diesem Artikel betrachtet haben.
